\begin{pastexam}[label={px:gk3-101-ds-02}]{高考三級資訊處理 101 年 第二題}
（申論；25 分）

下列的虛擬碼程式片段中，I 和 S 均為遞迴函式（recursive function），I 和 S 的參數 $A$ 是一個整數
陣列；I 和 S 的參數 $i$ 為不為負的整數，主要是做為陣列 $A$ 的索引（index）。假設陣列 $A$ 的元素
個數為 $n$，且其索引值為 0 到 $n-1$ 之間的數值。虛擬碼 swap $x$ and $y$ 的意思是將變數 $x$ 與變數
$y$ 的儲存值互換；亦即執行之後變數 $x$ 的儲存值為執行前變數 $y$ 的儲存值，執行之後變數 $y$ 的
儲存值為執行前變數 $x$ 的儲存值。令 $T(n)$ 為呼叫函式 I($A$, $n-1$) 的執行時間。$T(n)$ 會隨著
陣列 $A$ 所儲存的數值不同而有所不同。
\begin{quote}
S($A$, $i$) \{\\
\hspace*{2em}If $i$ <= 0, then return;\\
\hspace*{2em}S($A$, $i - 1$);\\
\hspace*{2em}I($A$, $i$);\\
\hspace*{2em}Return; \}\\
I($A$, $i$) \{\\
\hspace*{2em}If $i$ <= 0, then return;\\
\hspace*{2em}If $A[i] < A[i - 1]$ \{\\
\hspace*{4em}swap $A[i]$ and $A[i - 1]$ ;\\
\hspace*{4em}I($A$, $i - 1$); \}\\
\hspace*{2em}Return; \}
\end{quote}
\begin{itemize}[leftmargin=2em, labelsep=0.3em, itemsep=2pt]
\item[㈠] 請用 $O$-notation 表示 $T(n)$ 的上界（upper bound）；請用 $\Omega$-notation 表示 $T(n)$ 的
下界（lower bound）。（5 分）
\item[㈡] 請說明 $T(n)$ 最大時，程式開始執行前陣列 $A$ 所儲存的數值有何特性？理由為何？（5 分）
\item[㈢] 請問函式 S($A$, $n-1$) 的時間複雜度為何？請說明理由。（5 分）
\item[㈣] 請問執行函式 S($A$, $n-1$) 後，陣列 $A$ 儲存的內容有何特性？請證明你的觀察。（10 分）
\end{itemize}
\end{pastexam}

\begin{pastsoln}
I($A, i$) 把 $A[i]$ 一路往左換，直到左邊的數不比它大為止，正是插入排序「把一個數插進前面排好的
部分」那一步；S 則是遞迴版的插入排序（資結書第~1.3~節、虛擬碼~1.3）。 下面設 $n \ge 2$。

\textbf{㈠} 每次呼叫 I 只做常數個動作（一次比較、最多一次交換），再至少遞迴一次；遞迴時索引減一，
所以從 I($A, n-1$) 開始，呼叫鏈最長是 I($A, n-1$), I($A, n-2$), $\ldots$, I($A, 0$)，共 $n$ 次；
最短是第一次比較就發現 $A[n-1] \ge A[n-2]$，只有 $1$ 次呼叫、$1$ 次比較。 因此
\[
T(n) = O(n), \qquad T(n) = \Omega(1).
\]
兩個界都是緊的：最壞 $\Theta(n)$、最好 $\Theta(1)$，㈡給出達到最壞的輸入。

\textbf{㈡ $T(n)$ 最大的條件：$A[n-1]$ 比 $A[0], A[1], \ldots, A[n-2]$ 每一個都小}（嚴格小於），
其他位置的數怎麼排都沒關係。

理由：令 $x = A[n-1]$（原來的值）。 I($A, n-1$) 比較 $x$ 和 $A[n-2]$，若 $x < A[n-2]$ 就交換，$x$
到了第 $n-2$ 格，再呼叫 I($A, n-2$) 比較 $x$ 和 $A[n-3]$……。 每一次比較的另一邊 $A[j-1]$ 都還是
\textbf{原來的值}（只有 $x$ 在移動，左邊的數還沒被動過）。 所以 $x$ 會一路換到第 $0$ 格（$n - 1$
次比較、$n - 1$ 次交換、$n$ 次呼叫，工作量最多）當且僅當 $x < A[j]$ 對所有 $0 \le j \le n-2$ 都
成立，也就是 $x$ 是全陣列唯一的最小值；只要某個 $A[j] \le x$，$x$ 最晚停在第 $j+1$ 格，交換和呼叫
都至少少一次。 （只看比較次數的話，$x$ 比 $A[1..n-2]$ 都小時就已經比了 $n-1$ 次，但交換和呼叫
各少一次，執行時間仍以「$x$ 是唯一最小值」時最大。） 例如 $\langle 5, 9, 2, 7, 1 \rangle$：$1$
依序和 $7, 2, 9, 5$ 比較並交換，換到第 $0$ 格，達到最大值。

\textbf{㈢ S($A, n-1$) 的時間：最壞 $\Theta(n^2)$，最好 $\Theta(n)$；所以是 $O(n^2)$、$\Omega(n)$。}

S($A, i$) 先遞迴做 S($A, i-1$)，再做 I($A, i$)，所以 S($A, n-1$) 依序執行
I($A, 1$), I($A, 2$), $\ldots$, I($A, n-1$)，外加 $n$ 次 S 的呼叫（常數時間）：
\[
T_S(n) = \Theta(n) + \sum_{i=1}^{n-1} (\text{I}(A, i) \text{ 的時間}).
\]
I($A, i$) 的時間介於 $\Theta(1)$ 和 $\Theta(i)$ 之間（由㈠）。
\begin{itemize}[nosep]
\item 最壞：陣列一開始嚴格由大到小（例如 $\langle 5, 4, 3, 2, 1 \rangle$）。 由㈣，執行
I($A, i$) 前 $A[0..i-1]$ 已排好，而 $A[i]$ 比它們都小，由㈡每個 I($A, i$) 都換到底，比較 $i$
次，合計 $\sum_{i=1}^{n-1} i = n(n-1)/2$，$\Theta(n^2)$。
\item 最好：陣列一開始已由小到大排好，每個 I($A, i$) 比較一次就停，合計 $n - 1$ 次，$\Theta(n)$。
\end{itemize}
不論哪種輸入都不會超過 $n(n-1)/2$ 次比較，所以 $T_S(n) = O(n^2)$；最壞情況的時間複雜度是
$\Theta(n^2)$。 另外 S 的遞迴深度是 $n$，I 的遞迴深度最多 $n$，呼叫堆疊需要 $O(n)$ 的空間。

\textbf{㈣ 特性：執行 S($A, n-1$) 後，$A[0..n-1]$ 由小到大排好（非遞減），而且是原來那些數的
重新排列。}

重新排列這一點很明顯：程式只做交換。 下面證明排好。

\textbf{引理。} 若呼叫 I($A, i$) 前 $A[0..i-1]$ 已非遞減，則呼叫後 $A[0..i]$ 非遞減，且 $A[i+1..]$
不變。

\textbf{證明。} 對 $i$ 歸納。 $i = 0$：直接返回，$A[0..0]$ 只有一個數，成立。 設 $i \ge 1$，引理
對 $i - 1$ 成立。 令 $x = A[i]$、$m = A[i-1]$；因 $A[0..i-1]$ 非遞減，$m$ 是 $A[0..i-1]$ 的最大值。
\begin{itemize}[nosep]
\item 若 $x \ge m$：不交換，直接返回。 $A[0..i-1]$ 非遞減且 $A[i-1] = m \le x = A[i]$，所以
$A[0..i]$ 非遞減。
\item 若 $x < m$：交換後 $A[i] = m$、$A[i-1] = x$，$A[0..i-2]$ 沒動，仍非遞減。 由歸納假設，
I($A, i-1$) 之後 $A[0..i-1]$ 非遞減，而且它是 $\{A[0..i-2]$ 原來的數$\} \cup \{x\}$ 的重新排列，
這些數都 $\le m$（$A[0..i-2] \le m$，$x < m$）。 又 $A[i] = m$ 沒被動過，所以 $A[0..i]$ 非遞減。
\end{itemize}
I 只動到索引 $\le i$ 的格子，$A[i+1..]$ 不變。 \hfill$\blacksquare$

\textbf{定理。} 呼叫 S($A, i$) 後 $A[0..i]$ 非遞減。

\textbf{證明。} 對 $i$ 歸納。 $i \le 0$：直接返回，至多一個數，成立。 $i \ge 1$：由歸納假設，
S($A, i-1$) 之後 $A[0..i-1]$ 非遞減（且沒動到 $A[i]$ 以後）；由引理，接著的 I($A, i$) 使
$A[0..i]$ 非遞減。 \hfill$\blacksquare$

取 $i = n - 1$，整個陣列由小到大排好。 例：$\langle 3, 1, 2 \rangle$，S 先做 I($A, 1$)：$1 < 3$
交換得 $\langle 1, 3, 2 \rangle$；再做 I($A, 2$)：$2 < 3$ 交換得 $\langle 1, 2, 3 \rangle$，
再比 $2$ 和 $1$ 不換，結束。
\end{pastsoln}
